\section{Introduction}
\label{sec:intro}

\subsection{Motivation}
\label{sec:motivation}

World models that plan in a compact latent space enable fast
model-predictive control without decoding pixels at every step
\cite{ha2018worldmodels,hafner2019planet}.
Classical latent dynamics models often predict the next state with a
discrete transition (RNN, RSSM, or transformer block), then hand the
rollout to a planner such as CEM or MPPI.
Joint Embedding Predictive Architectures (JEPAs)
\cite{assran2023jepa} push the representation side further by
predicting future \emph{embeddings} rather than reconstructions,
avoiding a pixel decoder that can dominate compute and encourage
shortcut features.
Recent control-oriented JEPA stacks combine a ViT encoder trained from
scratch, an action-conditioned latent predictor, next-embedding MSE, and
SIGReg anti-collapse regularization~\cite{balestriero2025lejepa}, then
plan with CEM directly in latent space against a goal embedding.
Such systems can reach strong PushT success under full-data training,
but the predictor is typically still a discrete autoregressive block:
given context latents $z_t$ and action embeddings $a_t$, it directly
outputs $\hat z_{t+1}$.

Physical dynamics, however, are continuous in time.
A discrete next-step map must absorb the entire transition into a
single nonlinear layer, with no explicit notion of intermediate flow
along the planning interval.
That gap matters for temporal coherence, residual stability at
initialization, and---in principle---sharing structure across step
sizes when physical time is rescaled.
Recent continuous-time world models such as
\PriorODEWorld{}~\cite{liu2026odeworld} argue that latent velocity fields
$\mathrm{d}z/\mathrm{d}t = v_\theta(\cdot)$ better capture underlying
flow and support flexible temporal resolution.
Neural ODEs~\cite{chen2018neuralode} provide the standard tool for
parameterizing such fields and integrating them with classical
solvers.

We ask whether \emph{continuous} latent dynamics can serve as the
predictor inside an otherwise standard JEPA + SIGReg + CEM stack.
\method{} answers with an action-conditioned ODE velocity field and
RK4 integrate-to-next-step rule, reusing the same transformer capacity
as a discrete AR baseline under matched data and planning protocols.

\subsection{Contributions}
\label{sec:contributions}

\begin{enumerate}[leftmargin=*,itemsep=3pt]
\item \textbf{\method{} continuous latent predictor.}
We instantiate \method{} as an action-conditioned ODE inside a
JEPA-style visual world model:
$\mathrm{d}z/\mathrm{d}\tau = v_\theta(z,a)$, with one-step prediction
defined as classical RK4 integration over $\tau\in[0,1]$ (fixed
$N{=}4$ substeps), actions held constant on the interval, and a
zero-initialized linear velocity head for residual stability at
initialization.
The velocity backbone reuses the same ConditionalBlock / AdaLN-zero
transformer capacity as the discrete AR baseline
($\approx$18.07M parameters; +37{,}056 from the $v$-head).

\item \textbf{Matched 1k protocol on official data.}
We evaluate on the official Hugging Face PushT expert HDF5 with a
reproducible 1000-window split (seed~3072; 901 train / 99 val),
matching encoder, SIGReg ($\lambda{=}0.09$, $M{=}1024$), optimizer,
hardware batching (64$\times$accum~2), and CEM protocol between
discrete and ODE runs.
Indices are frozen and never regenerated between runs.

\item \textbf{Honest empirical findings.}
On this protocol, ODE lowers train/val loss relative to discrete AR
(0.234 / 0.506 vs.\ 0.248 / 0.536), while CEM success on the val-99
protocol is \textbf{tied} at 5/93~(5.38\%).
We explicitly contrast this with prior full-data PushT figures on the
order of $96\%$ success, and treat the 1k regime as a controlled
comparison rather than a claim of full-corpus performance.
No synthetic episodes are used.
\end{enumerate}

\paragraph{Why this matters.}
If continuous dynamics improve JEPA fit without destabilizing SIGReg or
CEM interfaces, visual world models can borrow ODE inductive bias without
abandoning latent MPC.
If they do not help planning under matched compute and data, that is
also useful: it suggests that predictor form is secondary to data scale
and planner configuration in the 1k regime.
Our results land in the middle---loss improves, CEM does not---and we
report both sides without overselling.

\paragraph{Scope.}
This paper is a controlled predictor comparison on PushT, not a
multi-environment bake-off or full-corpus training run.
Cube, Two-rooms, and Reacher remain future work
(\Cref{sec:limitations}).
We also do not convert or evaluate official Hugging Face pretrained
PushT weights under our val-99 protocol; that baseline remains future
work.
